\documentclass[14pt, a4paper]{article}
\usepackage{minitoc}
\usepackage[left=3.00cm, right=2.5cm, top=2.00cm, bottom=2.00cm]{geometry}
\usepackage{amsmath}
\usepackage{amssymb}
\usepackage{amsthm}
\usepackage{mathtools}
\usepackage{graphicx}
\usepackage{algpseudocode}
\usepackage{algorithm}
\usepackage{blindtext}
\usepackage{setspace}
\usepackage[utf8]{inputenc}
\usepackage[utf8]{vietnam}
\usepackage[center]{caption}
\usepackage[shortlabels]{enumitem}
\usepackage{fancyhdr} % header, footer
\usepackage{hyperref} % loại bỏ border với mục lục và công thức
\usepackage[nonumberlist, nopostdot, nogroupskip]{glossaries}
\usepackage{glossary-superragged}
\setglossarystyle{superraggedheaderborder}
\pagestyle{fancy}
%\usepackage[style=numeric,sortcites]{biblatex}
%\addbibresource{ref.bib}
%\usepackage[numbers]{natbib}
\usepackage{indentfirst}
\usepackage[natbib,backend=biber,style=ieee, sorting=ynt]{biblatex}
\bibliography{ref.bib}

\graphicspath{{./figures/}}

%\renewbibmacro*{cite}{%
%  \printtext[bibhyperref]{%
%    \printfield{prefixnumber}%
%    \printfield{labelnumber}%
%    \ifbool{bbx:subentry}%
%      {\printfield{entrysetcount}}%
%   \ifnumequal{\value{citecount}}{\value{citetotal}-1}%
%       {\gdef\multicitedelim{\addspace\bibstring{and}\space}}%
%       {\gdef\multicitedelim{\addcomma\space}}%
%    }%
%}

%\makenoidxglossaries
%
%% Danh mục thuật ngữ
%\newglossaryentry{GMRES}
%{
%	name={GMRES},
%	description={Generalized Minimal Residual Method}
%}
%
%\newglossaryentry{MINRES}
%{
%	name={MINRES},
%	description={Minimal Residual Method}
%}
%
%\hypersetup{
%    colorlinks=false,
%    pdfborder={0 0 0},
%}
%
%\title{Tiểu luận phương pháp số cho đại số tuyến tính}
%
%\author{Nguyễn Chí Thanh}
%
%%\date{24-04-2022}
\fancyhf{}
\rhead{\textbf{Môn học: Tối ưu hóa nâng cao}}
\lhead{\textbf{GVHD: TS. Hoàng Nam Dũng}}
\rfoot{\thepage}
\lfoot{\textbf{Học viên thực hiện: Nguyễn Chí Thanh - Nguyễn Đức Thịnh}}
\renewcommand{\headrulewidth}{0.4pt}
\renewcommand{\footrulewidth}{0.4pt}
%
%\numberwithin{equation}{section}
%\numberwithin{algorithm}{section}
%\numberwithin{figure}{section}
%
%\setlength{\parindent}{0.5cm}
%
%\setcounter{secnumdepth}{3} % Cho phép subsubsection trong report
%\setcounter{tocdepth}{3} % Chèn subsubsection vào bảng mục lục

%\newtheorem{dl}{Định lý}
%\newtheorem{md}{Mệnh đề}
%\newtheorem{bd}{Bổ đề}
%\newtheorem{dn}{Định nghĩa}
%\newtheorem{hq}{Hệ quả}

%\newtheorem{baitap}{Bài tập}
%\newtheorem*{loigiai}{Lời giải}

%\numberwithin{dl}{section}
%\numberwithin{md}{section}
%\numberwithin{bd}{section}
%\numberwithin{dn}{section}
%\numberwithin{hq}{section}

\setlength{\parindent}{0cm}

\newtheorem{dl}{Định lý}
\newtheoremstyle{sltheorem}
{}                % Space above
{}                % Space below
{\normalfont}        % Theorem body font % (default is "\upshape")
{}                % Indent amount
{\bfseries}       % Theorem head font % (default is \mdseries)
{.}               % Punctuation after theorem head % default: no punctuation
{ }               % Space after theorem head
{}                % Theorem head spec
\theoremstyle{sltheorem}
\newtheorem{baitap}{Bài tập}
\newtheoremstyle{soltheorem}
{}                % Space above
{}                % Space below
{\normalfont}        % Theorem body font % (default is "\upshape")
{}                % Indent amount
{\bfseries}       % Theorem head font % (default is \mdseries)
{.}               % Punctuation after theorem head % default: no punctuation
{\newline}               % Space after theorem head
{}                % Theorem head spec
\theoremstyle{soltheorem}
\newtheorem*{loigiai}{Lời giải}

\onehalfspacing

\begin{document}

    \begin{titlepage}

        \newcommand{\HRule}{\rule{\linewidth}{0.5mm}} % Defines a new command for the horizontal lines, change thickness here

        \center % Center everything on the page

        %----------------------------------------------------------------------------------------
        %	HEADING SECTIONS
        %----------------------------------------------------------------------------------------
        \textsc{\LARGE Đại học Quốc Gia Hà Nội}\\[0.5cm]
        \textsc{\LARGE Trường đại học Khoa học tự nhiên}\\[0.5cm] % Name of your university/college
        \textsc{\LARGE Khoa Toán - Cơ - Tin học}\\[0.5cm]

        \includegraphics[scale=0.2]{HUS-logo.jpg}\\[0.5cm]

        \textsc{\Large Chuyên ngành: Khoa học dữ liệu}\\[0.5cm] % Major heading such as course name


        %----------------------------------------------------------------------------------------
        %	TITLE SECTION
        %----------------------------------------------------------------------------------------

        \HRule \\[0.4cm]
        { \huge \bfseries Bài tập môn học}\\[0.4cm] % Title of your document
        \HRule \\[1.5cm]

        \textsc{\Large Môn học: Tối ưu hóa nâng cao}\\[1.5cm] % Minor heading such as course title


        \textsc{\Large Bài tập 1}\\[1.5cm]


        %----------------------------------------------------------------------------------------
        %	AUTHOR SECTION
        %----------------------------------------------------------------------------------------
        \begin{minipage}{0.4\textwidth}
            \begin{flushleft} \Large
            \emph{Giảng viên hướng dẫn:} \\
            TS. Hoàng Nam Dũng % Supervisor's Name
            \end{flushleft}
        \end{minipage}\\[1cm]

        \begin{minipage}{0.4\textwidth}
        \begin{flushleft} \Large
        \emph{Nhóm học viên thực hiện:}\\
        Nguyễn Chí Thanh \\
        MSHV: 21007925 \\ % Your name
        Nguyễn Đức Thịnh \\
        MSHV: 21007923 \\
        Lớp: Khoa học dữ liệu - K4
        \end{flushleft}
        \end{minipage}


        % If you don't want a supervisor, uncomment the two lines below and remove the section above
        %\Large \emph{Author:}\\
        %John \textsc{Smith}\\[3cm] % Your name

        %----------------------------------------------------------------------------------------
        %	DATE SECTION
        %----------------------------------------------------------------------------------------

        % I don't want day because it is English
        % {\large \today}\\[2cm] % Date, change the \today to a set date if you want to be precise

        %----------------------------------------------------------------------------------------
        %	LOGO SECTION
        %----------------------------------------------------------------------------------------

        %\includegraphics{logo/rsz_3logo-khtn.png}\\[1cm] % Include a department/university logo - this will require the graphicx package

        %----------------------------------------------------------------------------------------

        \vfill % Fill the rest of the page with whitespace

    \end{titlepage}

    \nocite{*}

    \newpage

    \begin{baitap}
        Hãy chứng minh các hàm sau là hàm lồi

        \begin{enumerate}[wide, labelwidth=!, labelindent=0pt,label=\textbf{\arabic*}.]
            \item Ridge Regression
            \begin{equation*}
                f(x)=\lVert Ax - b \rVert_2^2 + \lambda \lVert x \rVert_2^2
            \end{equation*}
            \item Lasso Regression
            \begin{equation*}
                f(x)=\lVert Ax - b \rVert_2^2 + \lambda \lVert x \rVert_1
            \end{equation*}
            \item Logistic Regression
            \begin{equation*}
                f(w)=\dfrac{1}{n}\sum_{i=1}^n \Big( -y_i x_i^T w + \ln(1 + e^{x_i^T w}) \Big)
            \end{equation*}
            với $(x_i, y_i) \in \mathbb{R}^{k+1},i=1,2,\dots,n$, là cặp dữ liệu đầu vào và nhãn tương ứng và $w \in \mathbb{R}^k$
        \end{enumerate}
    \end{baitap}

    \begin{loigiai}

        Hãy chứng minh các hàm sau là hàm lồi
        \begin{enumerate} [wide, labelwidth=!, labelindent=0pt,label=\textbf{\arabic*}.]
            \item Ridge Regression
            \begin{equation*}
                f(x)=\lVert Ax - b \rVert_2^2 + \lambda \lVert x \rVert_2^2
            \end{equation*}

            Ta đặt kích thước các ma trận $A \in \mathbb{R}^{m \times n}, b \in \mathbb{R}^{m}, x \in \mathbb{R}^{n}$, ta có $\mathrm{dom}(f)=\mathbb{R}^{n}$ là tập lồi

            \begin{equation*}
                \begin{aligned}
                    f(x) &= (Ax - b)^T (Ax - b) + \lambda x^T x \\
                    &=x^T A^T A x - x^TA^Tb - b^T Ax + b^T b + \lambda x^T x
                \end{aligned}
            \end{equation*}
            Do $x^TA^Tb$ và $b^TAx$ là hai số vô hướng và $(x^TA^Tb)^T=b^TAx$ nên $x^TA^Tb=b^TAx$ (chuyển vị của một số vô hướng bằng chính nó) nên:

            \begin{equation*}
                \begin{aligned}
                    f(x) &= x^T A^T A x - x^TA^Tb - b^T Ax + b^T b + \lambda x^T x\\
                    &=x^T A^T A x-2x^TA^Tb+b^Tb + \lambda x^T x
                \end{aligned}
            \end{equation*}
            Ta tính Gradient của $f(x)$ theo $x$:

            \begin{equation*}
                \nabla f(x) =\Big\lbrack A^TA + (A^TA)^T \Big\rbrack x - 2A^Tb + 2\lambda x = 2A^T A x - 2 A^T b + 2\lambda x
            \end{equation*}

            Ta tính ma trận Hessian của $f(x)$ theo $x$:

            \begin{equation*}
                \nabla^2 f(x) = 2A^T A + 2 \lambda I
            \end{equation*}

            Ta xét:

            \begin{equation*}
                p^T \nabla^2 f(x) p = 2 p^T A^T A p + 2\lambda p^Tp = 2 (Ap)^TAp + 2\lambda p^Tp \geq 0 \thickspace\forall \thickspace p \in \mathbb{R}^{n}
            \end{equation*}

            nên:

            \begin{equation*}
                \nabla^2 f(x) \succeq 0 \thickspace \forall x \in \mathbb{R}^{n}
            \end{equation*}

            Vì $\mathrm{dom}(f)$ là tập lồi và $\nabla^2 f(x) \succeq 0$ nên hàm $f(x)$ là một hàm lồi.

            \item Lasso Regression
            \begin{equation*}
                f(x)=\lVert Ax - b \rVert_2^2 + \lambda \lVert x \rVert_1
            \end{equation*}

            Ta đặt kích thước các ma trận $A \in \mathbb{R}^{m \times n}, b \in \mathbb{R}^{m}, x \in \mathbb{R}^{n}$, ta có $\mathrm{dom}(f)=\mathbb{R}^{n}$ là tập lồi,
            $\mathrm{dom}(\lVert Ax - b \rVert_2^2)=\mathbb{R}^n$ và $\mathrm{dom}(\lVert x \rVert_1)=\mathbb{R}^n$

            Ta chứng minh $\lVert Ax - b \rVert_2^2$ là hàm lồi:

            \begin{equation*}
                \begin{aligned}
                    \lVert Ax - b \rVert_2^2&=(Ax-b)^T(Ax-b)\\
                    &=x^TA^TAx - x^TA^Tb - b^TAx + b^Tb
                \end{aligned}
            \end{equation*}
            Do $x^TA^Tb$ và $b^TAx$ là hai số vô hướng và $(x^TA^Tb)^T=b^TAx$ nên $x^TA^Tb=b^TAx$ (chuyển vị của một số vô hướng bằng chính nó) nên:

            \begin{equation*}
                \lVert Ax - b \rVert_2^2=x^TA^TAx - 2x^TA^Tb + b^Tb
            \end{equation*}

            Ta tính Gradient của $\lVert Ax - b \rVert_2^2$ theo $x$:

            \begin{equation*}
                \nabla \lVert Ax - b \rVert_2^2=\Big \lbrack A^TA + (A^TA)^T \Big \rbrack x-2A^Tb=2A^TAx - 2A^Tb
            \end{equation*}

            Ta tính ma trận Hessian của $\lVert Ax - b \rVert_2^2$ theo $x$:
            \begin{equation*}
                \nabla^2 \lVert Ax - b \rVert_2^2=2A^TA
            \end{equation*}

            Ta xét:

            \begin{equation*}
                2p^TA^TAp=2(Ap)^TAp\geq0 \thickspace \forall p \in \mathbb{R}^{n} \Rightarrow 2A^TA \succeq 0 \thickspace \forall x \in \mathbb{R}^{n}
            \end{equation*}

            Mặt khác, $\mathrm{dom}(\lVert Ax - b \rVert_2^2)=\mathbb{R}^n$ là tập lồi. Vậy nên $\lVert Ax - b \rVert_2^2$ là hàm lồi

            Ta chứng minh $\lVert x \rVert_1$ là hàm lồi, với $\theta \in \lbrack 0, 1\rbrack, x,y \in \mathbb{R}^n$:

            \begin{equation*}
                \begin{aligned}
                    \lVert \theta x + (1-\theta)y \rVert_1 \leq \theta \lVert x \rVert_1 + (1-\theta)\lVert y \rVert_1 \thickspace \forall \thickspace \theta \in \lbrack 0, 1 \rbrack, x, y \in \mathbb{R}^n
                \end{aligned}
            \end{equation*}

            theo bất đẳng thức tam giác và $\mathrm{dom}(\lVert x \rVert_1)=\mathbb{R}^n$ là tập lồi nên $\lVert x \rVert_1$ là hàm lồi. 
            $\mathrm{dom}(f)$ là một tập lồi và tổng $\lVert Ax - b \rVert_2^2 + \lambda \lVert x \rVert_1 \thickspace (\lambda \geq 0)$ có các thành phần $\lVert Ax - b \rVert_2^2$ và $\lVert x \rVert_1$ lồi nên $f(x)$ là một hàm lồi.

            \item Logistic Regression
            \begin{equation*}
                f(w)=\dfrac{1}{n}\sum_{i=1}^n \Big( -y_i x_i^T w + \ln(1 + e^{x_i^T w}) \Big)
            \end{equation*}
            với $(x_i, y_i) \in \mathbb{R}^{k+1},i=1,2,\dots,n$ là cặp dữ liệu đầu vào và nhãn tương ứng và $w \in \mathbb{R}^k$

            Ta có $\mathrm{dom}(f(w)=\mathbb{R}^{k}$, là một tập lồi. Ta tính Gradient của $f(w)$ theo $w$:

            \begin{equation*}
                \nabla f(w)=\dfrac{1}{n}\sum_{i=1}^n \Big( -y_i x_i + \dfrac{e^{x_i^Tw}}{1 + e^{x_i^T w}}x_i \Big)
            \end{equation*}

            Ta tính ma trận Hessian của $f(w)$ theo $w$:

            \begin{equation*}
                \nabla^2 f(x) = \dfrac{1}{n}\sum_{i=1}^n \dfrac{e^{x_i^Tw}(1 + e^{x_i^T w})-e^{2x_i^T w}}{(1 + e^{x_i^T w})^2}x_i x_i^T=\dfrac{1}{n}\sum_{i=1}^n \dfrac{e^{x_i^Tw}}{(1 + e^{x_i^T w})^2}x_i x_i^T
            \end{equation*}

            Ta xét số hạng thứ $i$:

            \begin{equation*}
                \dfrac{e^{x_i^Tw}}{(1 + e^{x_i^T w})^2}x_i x_i^T
            \end{equation*}

            Ta xét ma trận Hessian $\nabla^2 f(x)$:

            \begin{equation*}
                \dfrac{e^{x_i^Tw}}{(1 + e^{x_i^T w})^2}p^Tx_i x_i^T p=\Bigg(\sqrt{\dfrac{e^{x_i^Tw}}{(1 + e^{x_i^T w})^2}} x_i^Tp\Bigg)^T\Bigg(\sqrt{\dfrac{e^{x_i^Tw}}{(1 + e^{x_i^T w})^2}} x_i^Tp\Bigg)\geq 0 \thickspace \forall p \in \mathbb{R}^{k} \thickspace \forall x \in \mathbb{R}^{k}
            \end{equation*}
            \begin{equation*}
                \begin{aligned}
                    \Rightarrow p^T \nabla^2 f(x)p&=\dfrac{1}{n}\sum_{i=1}^n \dfrac{e^{x_i^Tw}}{(1 + e^{x_i^T w})^2}p^Tx_i x_i^T\\
                    &=\dfrac{1}{n}\sum_{i=1}^n \Bigg(\sqrt{\dfrac{e^{x_i^Tw}}{(1 + e^{x_i^T w})^2}} x_i^Tp\Bigg)^T\Bigg(\sqrt{\dfrac{e^{x_i^Tw}}{(1 + e^{x_i^T w})^2}} x_i^Tp\Bigg)\geq 0 \thickspace \forall p \in \mathbb{R}^{k} \thickspace \forall x \in \mathbb{R}^{k}
                \end{aligned}
            \end{equation*}

            Ta có $\mathrm{dom}(f(w))=\mathbb{R}^{k}$ là một tập lồi và $\nabla^2 f(x) \succeq 0$ nên $f(w)$ là hàm lồi
        \end{enumerate}
        

    \end{loigiai}

    \begin{baitap}
        Cho hàm số $f(x, y) = x^4 + x^2 - 6xy + 3y^2$. Hãy tìm các điểm cực trị địa phương của $f$, hãy chỉ rõ đó là cực đại hay cực tiểu địa phương
    \end{baitap}

    \begin{loigiai}
        Ta có $\mathrm{dom}(f)=\mathbb{R}^2$ là một tập lồi. Ta tính Gradient của $f(x, y)$ theo $x$ và $y$:

        \begin{equation*}
            \nabla f(x, y) = \begin{bmatrix} \dfrac{\partial f(x, y)}{\partial x} \\ \dfrac{\partial f(x, y)}{\partial y}\end{bmatrix} = \begin{bmatrix} 4x^3 + 2x - 6y \\ -6x + 6y \end{bmatrix}
        \end{equation*}
        Ta tính ma trận Hessian của $f(x, y)$ theo $x$ và $y$:

        \begin{equation*}
            \nabla^2 f(x,y)=\begin{bmatrix} \dfrac{\partial^2 f(x, y)}{\partial x^2} & \dfrac{\partial^2 f(x, y)}{\partial x\partial y} \\ \dfrac{\partial^2 f(x, y)}{\partial y \partial x } & \dfrac{\partial^2 f(x, y)}{\partial y^2 } \end{bmatrix}=\begin{bmatrix} 12x^2 + 2 & -6 \\ -6 & 6 \end{bmatrix}
        \end{equation*}

        Ta nhận thấy $\nabla^2 f(x, y)$ tồn tại và liên tục với mọi $(x, y) \in \mathbb{R}^{2}$

        Để xét dấu của ma trận Hessian $\nabla^2 f(x, y)$, ta sẽ sử dụng một vài định lý và tiêu chuẩn.

        \begin{dl}[Tiêu chuẩn Sylvester] \label{dl:Sylvester}
            Cho ma trận $A \in \mathbb{R}^{n \times n}$ đối xứng, 
            đặt $A_k$ là ma trận thu được bằng cách lấy một ma trận con kích thước $k\times k$ góc trên bên trái của ma trận $A$. 
            Hơn nữa, ta đặt $\Delta_k=\det(A_k)$ là định thức con chính cấp $k$ của ma trận $A$.
            Điều kiện cần và đủ để ma trận $A$ là một ma trận xác định dương là:
            \begin{equation*}
                \Delta_k > 0 \thickspace \forall k = 1,\dots,n
            \end{equation*}
        \end{dl}

        Từ định lý \ref{dl:Sylvester}, để một ma trận $A \in \mathbb{R}^{n \times n}$ là một ma trận xác định âm thì điều kiện cần và đủ là:

        \begin{equation*}
            (-1)^k \Delta_k > 0 \thickspace \forall k = 1,\dots,n
        \end{equation*}

        Định lý \ref{dl:Sylvester} chỉ giúp ta xác định một ma trận đối xứng là ma trận xác định dương hay xác định âm mà không giúp ta xác định được một ma trận là nửa xác định âm hay nửa xác định dương.
        Liên quan đến ma trận nửa xác định dương ta có tiểu chuẩn Sylvester mở rộng.

        \begin{dl}[Tiêu chuẩn Sylvester mở rộng] \label{dl:Extended-Sylvester}
            Cho ma trận $A \in \mathbb{R}^{n \times n}$ đối xứng, 
            đặt $A_k$ là ma trận thu được bằng cách lấy một ma trận con kích thước $k\times k$ góc trên bên trái của ma trận $A$. 
            Hơn nữa, ta đặt $\Delta_k=\det(A_k)$ là định thức con chính cấp $k$ của ma trận $A$.
            Điều kiện cần để ma trận $A$ là một ma trận nửa xác định dương là:
            \begin{equation*}
                \Delta_k \geq 0 \thickspace \forall k = 1,\dots,n
            \end{equation*}
        \end{dl}

        Một điều cần chú ý rằng định lý \ref{dl:Extended-Sylvester} chỉ nói đến điều kiện cần để một ma trận đối xứng là một ma trận nửa xác định dương.
        Nếu một ma trận đối xứng nửa xác định dương thì sẽ có $\Delta_k \geq 0 \thickspace \forall k = 1,\dots,n$. 
        Nhưng nếu một ma trận đối xứng và $\Delta_k \geq 0 \thickspace \forall k = 1,\dots,n$ ta chưa thể kết luận rằng ma trận này là một ma trận nửa xác định dương.
        Nhưng nếu ta chỉ cần chỉ ra tồn tại ít nhất một $k \in \lbrack 1, n \rbrack$ sao cho $\Delta_k < 0$ thì ta có thể kết luận rằng ma trận này là ma trận không nửa xác định dương.

        Từ định lý \ref{dl:Extended-Sylvester}, nếu một ma trận đối xứng thì điều kiện cần để ma trận này nửa xác định âm là:

        \begin{equation*}
            (-1)^k \Delta_k \geq 0 \thickspace \forall k = 1,\dots,n
        \end{equation*}

        Với ma trận Hessian $\nabla^2 f(x, y)$ của bài toán, ma trận Hessian $\nabla^2 f(x, y)$ là ma trận đối xứng, ta xét các định thức con chính:
        Ta xét định thức con chính cấp 1:
        \begin{equation*}
            \det(\begin{bmatrix} 12x^2 + 2  \end{bmatrix})>0 \thickspace \forall (x, y) \in \mathbb{R}^2
        \end{equation*}

        Theo định lý \ref{dl:Extended-Sylvester}, để ma trận $\nabla^2 f(x, y)$ có cơ hội là ma trận nửa xác định âm thì $(-1)^k \Delta_k \geq 0 \thickspace \forall k = 1,\dots,n$ nhưng định thức con chính đầu tiên lớn hơn 0 $\forall (x, y) \in \mathbb{R}^2$ nên ma trận Hessian $\nabla^2 f(x, y)$ không thỏa mãn điều kiện cần để là một ma trận nửa xác định âm.
        Vậy ma trận Hessian $\nabla^2 f(x, y)$ không là một ma trận nửa xác định âm.
        Ta xét định thức con chính cấp 2:
        \begin{equation*}
            \det(\nabla^2 f(x,y))=\det\Bigg(\begin{bmatrix} 12x^2 + 2 & -6 \\ -6 & 6 \end{bmatrix} \Bigg)=6(12x^2 + 2) - 36 =72x^2 - 24
        \end{equation*}

        $\det\Big(\nabla^2 f(x,y)\Big) < 0$ với $\forall x \in \Big(-\dfrac{1}{\sqrt{3}}; \dfrac{1}{\sqrt{3}}\Big)$, 
        ma trận Hessian $\nabla^2 f(x,y)$ không thỏa mãn điều kiện cần để là một ma trận nửa xác định dương do $\forall x \in \Big(-\dfrac{1}{\sqrt{3}}; \dfrac{1}{\sqrt{3}}\Big)$ định thức con chính cấp 2 của $\nabla^2 f(x, y)$ nhỏ hơn 0.
        Như vậy  ma trận Hessian $\nabla^2 f(x,y)$ không là một ma trận nửa xác định dương $\forall (x, y) \in \mathbb{R}^{2}$.


        Ma trận Hessian $\nabla^2 f(x,y)$ vừa là một ma trận không nửa xác định dương vừa là một ma trận không nửa xác định âm nên ma trận $\nabla^2 f(x,y)$ là ma trận không xác định dấu. Hàm $f(x, y)$ không là một hàm lồi (cũng không phải là một hàm lõm).

        Để tìm các cực trị địa phương của hàm $f(x, y)$, ta cần tìm tất cả các điểm dừng và đánh giá dấu của ma trận Hessian $\nabla^2 f(x,y)$ tại các điểm dừng này.
        
        Để tìm các điểm dừng, ta giải phương trình $\nabla f(x, y) = 0$

        \begin{equation*}
            \nabla f(x, y) = \begin{bmatrix} 4x^3 + 2x - 6y \\ -6x + 6y \end{bmatrix} = \begin{bmatrix} 0 \\ 0 \end{bmatrix}
        \end{equation*}

        \begin{equation}
            \begin{aligned}
            &\Leftrightarrow \begin{cases} 4x^3 + 2x - 6y = 0 \\ x = y \end{cases} \\
            &\Leftrightarrow \begin{cases} 4x(x^2-1) = 0 \\ x = y \end{cases} \\
            &\Leftrightarrow \left [\begin{array}{l} x=y=0 \\ x=y=1 \\ x=y=-1 \end{array}\right.
            \end{aligned}
        \end{equation}

        Ta sẽ xét ma trận Hessian $\nabla^2 f(x,y)$ tại các điểm dừng là nghiệm của phương trình $\nabla f(x, y) = 0$:

        \begin{itemize}
            \item $\nabla^2 f(x=0,y=0)=\begin{bmatrix} 2 & -6 \\ -6 & 6 \end{bmatrix}$. 
            Với vector $p_1 = \begin{bmatrix} 0 \\ 1\end{bmatrix}$, ta có $p_1^T \nabla^2 f(x=0,y=0) p_1=\begin{bmatrix} 0 & 1\end{bmatrix}\begin{bmatrix} 2 & -6 \\ -6 & 6 \end{bmatrix} \begin{bmatrix} 0 \\ 1\end{bmatrix}=6 > 0$. 
            Vậy ma trận $\nabla^2 f(x=0,y=0)=\begin{bmatrix} 2 & -6 \\ -6 & 6 \end{bmatrix}$ là ma trận không nửa xác định âm. 
            Với $p_2 = \begin{bmatrix} 1 \\ 1\end{bmatrix}$, ta có $p_2^T \nabla^2 f(x=0,y=0) p_2=\begin{bmatrix} 1 & 1\end{bmatrix}\begin{bmatrix} 2 & -6 \\ -6 & 6 \end{bmatrix} \begin{bmatrix} 1 \\ 1\end{bmatrix}=-4 < 0$.
            Vậy ma trận $\nabla^2 f(x=0,y=0)=\begin{bmatrix} 2 & -6 \\ -6 & 6 \end{bmatrix}$ là ma trận không nửa xác định dương.
            Vì ma trận $\nabla^2 f(x=0,y=0)=\begin{bmatrix} 2 & -6 \\ -6 & 6 \end{bmatrix}$ đồng thời không nửa xác định âm và không nửa xác định dương nên $\nabla^2 f(x=0,y=0)$ là ma trận không xác định dấu. 
            Nên $(x, y)=(0, 0)$ không phải là một cực trị địa phương.
            \item $\nabla^2 f(x=1,y=1)=\begin{bmatrix} 14 & -6 \\ -6 & 6 \end{bmatrix}$ là một ma trận đối xứng. $\det{(\begin{bmatrix}14\end{bmatrix})}=14>0, \det{\Big(\begin{bmatrix} 14 & -6 \\ -6 & 6 \end{bmatrix}\Big)}=14\times6 - (-6)\times(-6)=48>0$.
            Theo định lý \ref{dl:Sylvester}, ma trận $\nabla^2 f(x=1,y=1)$ đối xứng và tất cả các định thức con chính đều lớn hơn 0, nên $\nabla^2 f(x=1,y=1)$ là một ma trận xác định dương.
            Vì $\nabla^2 f(x, y)$ tồn tại và liên tục với mọi $(x, y) \in \mathbb{R}^2$ nên $\nabla^2 f(x, y)$ cũng tồn tại là liên tục trong một lân cận mở của $(x,y)=(1,1)$, $\nabla f(x=1,y=1)=0$ và $\nabla^2 f(x=1,y=1)$ là ma trận xác định dương nên $(x,y)=(1,1)$ là một cực tiểu địa phương ngặt.
            \item $\nabla^2 f(x=-1,y=-1)=\begin{bmatrix} 14 & -6 \\ -6 & 6 \end{bmatrix}$ là một ma trận đối xứng. $\det{(\begin{bmatrix}14\end{bmatrix})}=14>0, \det{\Big(\begin{bmatrix} 14 & -6 \\ -6 & 6 \end{bmatrix}\Big)}=14\times6 - (-6)\times(-6)=48>0$.
            Theo định lý \ref{dl:Sylvester}, ma trận $\nabla^2 f(x=-1,y=-1)$ đối xứng và tất cả các định thức con chính đều lớn hơn 0, nên $\nabla^2 f(x=-1,y=-1)$ là một ma trận xác định dương.
            Vì $\nabla^2 f(x, y)$ tồn tại và liên tục với mọi $(x, y) \in \mathbb{R}^2$ nên $\nabla^2 f(x, y)$ cũng tồn tại là liên tục trong một lân cận mở của $(x,y)=(-1,-1)$, $\nabla f(x=-1,y=-1)=0$ và $\nabla^2 f(x=-1,y=-1)$ là ma trận xác định dương nên Vậy $x=y=-1$ là một cực tiểu địa phương ngặt.
        \end{itemize}
        Vậy ta có hai điểm $(x,y)=(1, 1)$ và $(x,y)=(-1,-1)$ là các điểm cực tiểu địa phương.
    \end{loigiai}

    \begin{baitap}
        Cho $A \in \mathbb{R}^{m \times n}$ với $m > n$ và $\mathrm{rank}(A)=n$, $b \in \mathbb{R}^{m}$. Hãy giải bài toán tối ưu sau:
        \begin{equation*}
            \min_x \lVert Ax - b \rVert_2^2
        \end{equation*}
    \end{baitap}

    \begin{loigiai}
        Ta có $\mathrm{dom}\Big( \lVert Ax-b \rVert_2^2\Big)=\mathbb{R}^{n}$ là một tập lồi. Ta sẽ chứng minh
        $\lVert Ax - b \rVert_2^2$ là hàm lồi. Ta xét:

        \begin{equation*}
            \begin{aligned}
            \lVert Ax - b \rVert_2^2&=(Ax-b)^T(Ax-b)\\
            &=x^TA^TAx - x^TA^Tb - b^TAx + b^Tb
            \end{aligned}
        \end{equation*}

        Do $x^TA^Tb$ và $b^TAx$ là hai số vô hướng và $(x^TA^Tb)^T=b^TAx$ nên $x^TA^Tb=b^TAx$ (chuyển vị của một số vô hướng bằng chính nó) nên:
        \begin{equation*}
            \lVert Ax - b \rVert_2^2=x^TA^TAx - 2x^TA^Tb + b^Tb
        \end{equation*}
        Ta nhận thấy $\lVert Ax - b \rVert_2^2$ là hàm khả vi tại mọi điểm $x \in \mathbb{R}^n$. Ta tính Gradient của $\lVert Ax - b \rVert_2^2$ theo $x$:

        \begin{equation*}
            \nabla \lVert Ax - b \rVert_2^2=\Big\lbrack A^TA + (A^TA)^T \Big\rbrack x-2A^Tb=2A^TAx - 2A^Tb
        \end{equation*}

        Ta tính ma trận Hessian của $\lVert Ax - b \rVert_2^2$ theo $x$:

        \begin{equation*}
            \nabla^2 \lVert Ax - b \rVert_2^2=2A^TA
        \end{equation*}

        Ta xét:

        \begin{equation*}
            2p^TA^TAp=2(Ap)^TAp\geq0 \thickspace \forall p \in \mathbb{R}^{n} \Rightarrow 2A^TA \succeq 0 \Rightarrow \nabla^2 \lVert Ax - b \rVert_2^2 \succeq 0
        \end{equation*}

        Ta có $\mathrm{dom}\Big( \lVert Ax-b \rVert_2^2\Big)=\mathbb{R}^{n}$ là một tập lồi và $\nabla^2 \lVert Ax - b \rVert_2^2 \succeq 0$. Vì vậy $\lVert Ax - b \rVert_2^2$ là hàm lồi. Mặt khác, $\lVert Ax - b \rVert_2^2$ khả vi tại mọi điểm $x \in \mathbb{R}^{n}$ nên mọi điểm dừng của hàm $\lVert Ax - b \rVert_2^2$ là cực tiểu toàn cục.

        Ta giải phương trình $\nabla \lVert Ax - b \rVert_2^2=0$:

        \begin{equation*}
            \begin{aligned}
                &\nabla \lVert Ax - b \rVert_2^2 = 2 A^T A x - 2 A^T b = 0 \\
                \Leftrightarrow &A^T A x=A^T b
            \end{aligned}
        \end{equation*}

        Với $\mathrm{rank}(A)=n$, ta chứng minh $A^T A$ khả nghịch.

        Ta xét hệ phương trình:

        \begin{equation*}
            A^T A x = 0 \Rightarrow x^T A^T A x = 0 \Leftrightarrow (Ax)^T (Ax) = 0 \Leftrightarrow \lVert Ax \rVert_2^2 =0 \Leftrightarrow Ax = 0
        \end{equation*}

        Do $\mathrm{rank}(A)=n$ nên các cột của $A$ độc lập tuyến tính vậy để $Ax=0 \Leftrightarrow x = 0$
        Vậy $A^T A x = 0 \Leftrightarrow x = 0$, nên các cột của $A^T A$ độc lập tuyến tính, ta suy ra $A^T A$ khả nghịch.

        Vậy ta có nghiệm tối ưu toàn cục của bài toán là nghiệm của phương trình $\nabla \lVert Ax - b \rVert_2^2=0$:

        \begin{equation*}
            x^* = \underset{x}{\mathrm{argmin}}\lVert Ax - b \rVert_2^2=(A^TA)^{-1}A^T b
        \end{equation*}
    \end{loigiai}

    \begin{baitap}
        Xét hàm số Rosenbrock
        \begin{equation*}
            f(x)=100(x_2 - x_1^2)^2 + (1-x_1)^2
        \end{equation*}
        \begin{enumerate}[wide, labelwidth=!, labelindent=0pt,label=\textbf{\arabic*}.]
            \item $f$ có phải là hàm lồi không? Tại sao?
            \item Hãy tìm tất cả các điểm cực tiểu địa phương của $f$. Hàm $f$ có điểm cực tiểu toàn cục không? Nếu có hãy chỉ ra điểm đó là điểm nào?
        \end{enumerate}
    \end{baitap}

    \begin{loigiai}
        Xét hàm số Rosenbrock
        \begin{equation*}
            f(x)=100(x_2 - x_1^2)^2 + (1-x_1)^2
        \end{equation*}
        \begin{enumerate} [wide, labelwidth=!, labelindent=0pt,label=\textbf{\arabic*}.]
            \item $f$ có phải là hàm lồi không? Tại sao?

            Ta có $\mathrm{dom}(f)=\mathbb{R}^{2}$ là một tập lồi. Ta tính Gradient của $f(x)$ theo $x_1, x_2$:

            \begin{equation*}
                \nabla f(x) = \begin{bmatrix} \dfrac{\partial f(x)}{\partial x_1} \\ \dfrac{\partial f(x)}{\partial x_2}\end{bmatrix} = \begin{bmatrix} -400x_1(x_2 - x_1^2) -2(1-x_1) \\ 200(x_2 - x_1^2) \end{bmatrix}
            \end{equation*}

            Ta tính ma trận Hessian của $f(x)$ theo $x_1, x_2$:

            \begin{equation*}
                \nabla^2 f(x)=\begin{bmatrix} \dfrac{\partial^2 f(x)}{\partial x_1^2} & \dfrac{\partial^2 f(x)}{\partial x_1\partial x_2} \\ \dfrac{\partial^2 f(x)}{\partial x_2 \partial x_1 } & \dfrac{\partial^2 f(x)}{\partial x_2^2 } \end{bmatrix} = \begin{bmatrix} -400 x_2 + 1200x_1^2 + 2 & -400 x_1 \\ -400x_1 & 200 \end{bmatrix}
            \end{equation*}

            Ta nhận thấy $\nabla^2 f(x)$ tồn tại và liên tục với mọi $(x_1, x_2) \in \mathbb{R}^{2}$

            Ma trận $\nabla^2 f(x)$ là ma trận đối xứng. Ta xét định thức con chính cấp 1 của $\nabla^2 f(x)$:

            \begin{equation*}
                \det (\begin{bmatrix} -400 x_2 + 1200x_1^2 + 2 \end{bmatrix})=-400 x_2 + 1200x_1^2 + 2
            \end{equation*}

            Dấu của định thức này phụ thuộc vào $x_1, x_2$.
            \begin{itemize}
                \item Nếu $x_2 > \dfrac{1200x_1^2 + 2}{400}$ thì định thức $\det (\begin{bmatrix} -400 x_2 + 1200x_1^2 + 2 \end{bmatrix})<0$. Theo định lý \ref{dl:Extended-Sylvester}, ma trận $\nabla^2 f(x)$ không đạt điều kiện cần để là một ma trận nửa xác định dương. Vì vậy $\nabla^2 f(x)$ không là một ma trận nửa xác định dương $\forall (x_1, x_2) \in \mathbb{R}^2$.
                \item Nếu $x_2 < \dfrac{1200x_1^2 + 2}{400}$ thì $\det (\begin{bmatrix} -400 x_2 + 1200x_1^2 + 2 \end{bmatrix})>0$. Theo định lý \ref{dl:Extended-Sylvester}, ma trận $\nabla^2 f(x)$ không đạt điều kiện cần để là một ma trận nửa xác định âm. Vì vậy $\nabla^2 f(x)$ không là một ma trận nửa xác định âm $\forall (x_1, x_2) \in \mathbb{R}^2$.
            \end{itemize}
            Ma trận Hessian $\nabla^2 f(x)$ vừa là một ma trận không nửa xác định dương và là một ma trận không nửa xác định âm nên $\nabla^2 f(x)$ là ma trận không xác định dấu.

            $\nabla^2 f(x)$ là ma trận không xác định dấu nên $f(x)$ không là hàm lồi.

            \item Hãy tìm tất cả các điểm cực tiểu địa phương của $f$. Hàm $f$ có điểm cực tiểu toàn cục không? Nếu có hãy chỉ ra điểm đó là điểm nào?
            
            Để tìm các điểm dừng, ta giải phương trình $\nabla f(x)=0$:

            \begin{equation*}
                \nabla f(x) = \begin{bmatrix} \dfrac{\partial f(x)}{\partial x_1} \\ \dfrac{\partial f(x)}{\partial x_2}\end{bmatrix} = \begin{bmatrix} -400x_1(x_2 - x_1^2) -2(1-x_1) \\ 200(x_2 - x_1^2) \end{bmatrix} = \begin{bmatrix} 0 \\ 0 \end{bmatrix}
            \end{equation*}

            \begin{equation*}
                \begin{aligned}
                    &\Leftrightarrow\begin{cases} -400x_1(x_2 - x_1^2) -2(1-x_1) = 0 \\ 200(x_2 - x_1^2) = 0 \end{cases} \\
                    &\Leftrightarrow \begin{cases} -400x_1(x_2 - x_1^2) -2(1-x_1) = 0 \\ x_2 = x_1^2 \end{cases} \\
                    &\Leftrightarrow \begin{cases} x_1=1 \\ x_2 = x_1^2=1 \end{cases} \\
                    &\Leftrightarrow \begin{cases} x_1=1 \\ x_2=1 \end{cases}
                \end{aligned}
            \end{equation*}
            Ta sẽ xét ma trận Hessian $\nabla^2 f(x)$ tại các điểm dừng nghiệm của phương trình $\nabla f(x) = 0$:

            \begin{equation*}
                \nabla^2 f(x_1=1, x_2=1)=\begin{bmatrix} 802 & -400 \\ -400 & 200 \end{bmatrix}
            \end{equation*}

            Ta có ma trận $\nabla^2 f(x_1=1, x_2=1)$ là ma trận đối xứng. Ta xét các định thức con chính cấp 1 và cấp 2 của ma trận $\nabla^2 f(x_1=1, x_2=1)$:

            \begin{equation*}
                \begin{aligned}
                    &\det(\lbrack 802 \rbrack)=802>0 \\
                    &\det{\Big(\begin{bmatrix} 802 & -400 \\ -400 & 200 \end{bmatrix} \Big)}=802\times200 - (-400)\times(-400)=400>0 \\
                    \Rightarrow&\nabla^2 f(x_1=1, x_2=1) \succ 0
                \end{aligned}
            \end{equation*}

            Theo định lý \ref{dl:Sylvester}, ma trận $\nabla^2 f(x_1=1, x_2=1)$ là ma trận xác định dương do $\nabla^2 f(x_1=1, x_2=1)$ đối xứng và tất cả các định thức con chính (cấp 1 và cấp 2) đều lơn 0.

            Vì $\nabla^2 f(x)$ tồn tại và liên tục với mọi $(x_1, x_2) \in \mathbb{R}^{2}$ nên $\nabla^2 f(x, y)$ cũng tồn tại và liên tục trong một lân cận mở của $(x_1, x_2)=(1,1)$, $\nabla f(x_1=1, x_2=1)=0$ và $\nabla^2 f(x_1=1, x_2=1)$ là ma trận xác định dương nên $(x_1, x_2)=(1, 1)$ là một cực tiểu địa phương ngặt. Ta có $f(x_1=1, x_2=1)=0$. Mặt khác:

            \begin{equation*}
                \begin{aligned}
                    &f(x)=100(x_2 - x_1^2)^2 + (1-x_1)^2 \geq 0 = f(x_1=1,x_2=1) \thickspace \forall (x_1, x_2) \in \mathbb{R}^{2} \\
                    \Rightarrow&f(x_1=1,x_2=1) \leq f(x_1, x_2) \thickspace \forall (x_1, x_2) \in \mathbb{R}^{2}
                \end{aligned}
            \end{equation*}
            Vì vậy $(x_1, x_2)=(1, 1)$ cũng là một cực tiểu toàn cục của hàm $f(x)$.

            Kết luận: Hàm $f(x)$ chỉ có một cực tiểu địa phương $(x_1, x_2)=(1,1)$ và cũng chính là cực tiểu toàn cục của hàm $f$.
        \end{enumerate}
    \end{loigiai}

    \newpage
    \printbibliography[title={TÀI LIỆU THAM KHẢO}]
\end{document}